\addchap{Zusammenfassung}
\begingroup
\emergencystretch=2em

Axionen und axionartige Teilchen (ALPs) sind gut motivierte Kandidaten für dunkle Materie. Ihre Gradientenkopplung an Kernspins erzeugt ein oszillierendes pseudomagnetisches Feld, das mittels Kernspinresonanz (NMR) nachweisbar ist.
In dieser Dissertation werden Signalmodelle, Frequenzscanstrategien und Datenanalysemethoden entwickelt und auf Suchen zum Machbarkeitsnachweis mit dem Cosmic Axion Spin Precession Experiment (CASPEr) angewendet.

Betrachtet werden zwei Szenarien: ein virialisierter Milchstraßenhalo mit stochastischen Signalen endlicher Kohärenzzeit und ein Halo gravitativ erdgebundener ALPs mit möglicherweise längeren Signalkohärenzzeiten und erhöhter lokaler Dichte.
Die Frequenzscananalyse zeigt, dass sich die Kopplungsempfindlichkeit bei langer Messdauer mit der gesamten Integrationszeit lediglich gemäß $T_\mathrm{tot}^{-1/4}$ verbessert. Längere Kernspin-Kohärenzzeiten und an die empfindliche Bandbreite angepasste Frequenzschritte sind daher für effiziente Suchen entscheidend.

Die CASPEr-Gradient-LowField-Apparatur kombiniert eine thermisch polarisierte Flüssigmethanolprobe mit SQUID-Magnetometrie.
Bei Suchen nahe \SI{1.35}{\MHz}, entsprechend ALP-Massen nahe \SI{5.58}{\neV/c^2}, wurden keine statistisch signifikanten Kandidaten gefunden.
Die Suche im Milchstraßenhalo liefert über ein Frequenzintervall von \SI{240}{\Hz} eine obere Grenze für die ALP-Proton-Kopplung von etwa $|g_\mathrm{ap}|\lesssim\SI{1e-2}{\GeV^{-1}}$ bei einem Konfidenzniveau von \SI{95}{\percent}.
Die Suche nach erdgebundenen ALPs liefert über \SI{48}{\Hz} eine ungefähre Grenze von $|g_\mathrm{ap}|\lesssim\SI{1e-9}{\GeV^{-1}}$ bei einem Konfidenzniveau von \SI{95}{\percent}, unter den Annahmen zur Dichte und Zustandsbesetzung des erdgebundenen Halos.

Diese Messungen demonstrieren die Flüssigkeits-NMR zur Suche nach der ALP-Proton-Gradientenkopplung unter verschiedenen Haloannahmen.
Die geringe thermische Protonenpolarisation von etwa \num{1.8e-7} begrenzt die Empfindlichkeit maßgeblich. Hyperpolarisation ermöglicht potenziell Verbesserungen um mehrere Größenordnungen, ergänzt durch geringeres Detektorrauschen und höhere Feldhomogenität.
Die entwickelten Signalmodelle, statistischen Methoden und numerischen Simulationen der Spindynamik bilden eine Grundlage für zukünftige Suchen mit stärker polarisierten Proben.
\par\endgroup
